We use the conventions and results of S. Boucksom, C. Favre, and M. Jonsson,
``Valuations and plurisubharmonic singularities,'' Publ. Res. Inst. Math.
Sci. 44 (2008), 449--494, arXiv:math/0702487v2, abbreviated BFJ.  The
relevant exact source facts are the following.

(i) BFJ Section 1.2 defines, for every projective smooth modification
\(\pi:X_\pi\to X\) that is an isomorphism away from (0\), the real vector
space \(\operatorname{Div}(\pi)\) of exceptional divisors.  It states that
an exceptional divisor is uniquely determined by its numerical class and
identifies \(\operatorname{Div}(\pi)\) with
\(NS(X_\pi/X)_{\mathbf R}\).  BFJ Proposition 1.14 identifies exceptional
Cartier b-divisors with piecewise-affine functions and states that the
integral Cartier group is generated by \(Z(\mathfrak a)\) as
\(\mathfrak a\) ranges over the primary ideals.  BFJ also states in Section
1.2 that (Z(\mathfrak a)) is Cartier exactly when \(\mathfrak a\) is
primary, and, on a model dominating the blowup of \(\mathfrak a\), its
incarnation is the negative of the effective divisor defined by the
principalized ideal.

(ii) BFJ Theorem 5.1 associates to a psh germ (w\) a nef Weil b-divisor
\(Z(w)\); its corresponding formal psh function is the valuative transform
\(\widehat w\).  BFJ Proposition 5.6 identifies its values on
quasimonomial valuations with the appropriate Kiselman numbers.  On a
divisorial valuation \(\nu_E=\operatorname{ord}_E/
\operatorname{ord}_E(\mathfrak m)\), one has
\(\widehat w(\nu_E)=-\nu_E(w)\).

(iii) BFJ Lemma 5.10 states that a psh weight whose exponential is Holder
continuous is tame.  BFJ Theorem B states that for every tame psh weight
\(\varphi\) and every psh germ (w\),
\[
 \nu_\varphi(w)=\int_{\mathcal V}-\widehat w\,
                    \operatorname{MA}(\widehat\varphi).
\tag{1}
\]
BFJ Propositions 4.6 and 4.9 identify this integral, when the
Monge--Ampere inputs are Cartier/piecewise affine, with the corresponding
b-divisor intersection, computed on a model determining those inputs by
using the trace of the remaining nef Weil b-divisor.  Explicitly, if
\(C_1,\ldots,C_{n-1}\) are nef Cartier b-divisors determined on \(X_\pi\),
then
\[
 \int_{\mathcal V}g_W\,
   \operatorname{MA}(g_{C_1},\ldots,g_{C_{n-1}})
 =W_\pi\cdot(C_1)_\pi\cdots(C_{n-1})_\pi
\tag{2}
\]
for a nef Weil b-divisor (W\).  Formula (2) follows directly as well from
the atomic formula in BFJ Proposition 4.9: the mass at
\(\nu_E\) is
\(\operatorname{ord}_E(\mathfrak m)
 E\cdot(C_1)_\pi\cdots(C_{n-1})_\pi\), while
\(g_W(\nu_E)=\operatorname{ord}_E(W)/
\operatorname{ord}_E(\mathfrak m)\).

We now prove the statement.

For an \(\mathfrak m\)-primary ideal
\(\mathfrak a=(f_1,\ldots,f_N)\), put
\[
                    \varphi_{\mathfrak a}=\log\max_j|f_j|.
\tag{3}
\]
It has an isolated pole, and
\(\exp(\varphi_{\mathfrak a})=\max_j|f_j|\) is locally Lipschitz.  Thus
\(\varphi_{\mathfrak a}\) is a tame psh weight by BFJ Lemma 5.10.  Its
valuative transform and divisor are
\[
 \widehat\varphi_{\mathfrak a}=\log|\mathfrak a|,
 \qquad Z(\varphi_{\mathfrak a})=Z(\mathfrak a).
\tag{4}
\]
For primary ideals \(\mathfrak a_i\) and positive integers (k_i\), the
product is again primary and
\[
 Z\!\left(\prod_i\mathfrak a_i^{k_i}\right)
                         =\sum_i k_iZ(\mathfrak a_i).
\tag{5}
\]

We next make precise the intersection functional that will be polarized.
If (W\) is a nef Weil b-divisor and
\(C_1,\ldots,C_{n-1}\) are arbitrary exceptional Cartier b-divisors, choose
a common determination \(\pi\) of the (C_i\) and set
\[
 I_W(C_1,\ldots,C_{n-1})
 :=W_\pi\cdot(C_1)_\pi\cdots(C_{n-1})_\pi.
\tag{6}
\]
This does not depend on the determination.  Indeed, if
\(\mu:X_{\pi'}\to X_\pi\) is a higher model, then
\(\mu_*W_{\pi'}=W_\pi\) and the (C_i\) pull back, so the projection
formula gives (6) on (X_\pi\).  Hence (I_W\) is a symmetric multilinear
functional of its Cartier arguments.  When those arguments are nef, (6)
is exactly the BFJ intersection/integral in (2).  Notice that there is only
one possibly non-Cartier Weil input (W\); no intersection of several
signed differences of nef Weil b-divisors is being introduced.

Define the symmetric multilinear form on exceptional Cartier b-divisors
\[
 \Lambda(C_1,\ldots,C_{n-1})
 :=I_{Z(u)}(C_1,\ldots,C_{n-1})
  -I_{Z(v)}(C_1,\ldots,C_{n-1}).
\tag{7}
\]
By the assumed equality, (1), (2), and (4),
\[
 \Lambda(Z(\mathfrak a),\ldots,Z(\mathfrak a))=0
\tag{8}
\]
for every primary ideal \(\mathfrak a\).

Fix primary ideals \(\mathfrak a_1,\ldots,\mathfrak a_s\), set
\(W_i=Z(\mathfrak a_i)\), and consider
\[
 P(t_1,\ldots,t_s)
 :=\Lambda\!\left(\sum_i t_iW_i,\ldots,\sum_i t_iW_i\right).
\tag{9}
\]
This is a homogeneous real polynomial of degree (n-1\).  By (5) and (8),
\(P(k_1,\ldots,k_s)=0\) for all positive integers (k_i\).  A polynomial
vanishing on \(\mathbf Z_{>0}^s\) is identically zero, by induction on (s\)
using the one-variable fact.  Every coefficient of (9) therefore vanishes,
and hence
\[
 \Lambda(W_{i_1},\ldots,W_{i_{n-1}})=0
\tag{10}
\]
for every choice of indices.  BFJ Proposition 1.14 says that the integral
Cartier b-divisor group is generated by the (Z(\mathfrak a)\) for primary
ideals.  Tensoring with \(\mathbf R\) and using multilinearity in (10) gives
\[
             \Lambda(C_1,\ldots,C_{n-1})=0
\tag{11}
\]
for all real exceptional Cartier b-divisors (C_i\).

We record the finite-model separation needed to use (11).  Let
\(\pi:Y\to X\) be one BFJ modification, where (n\ge2\), with exceptional
prime divisors (E_1,\ldots,E_r\).  Since \(\pi\) is projective, it has a
relatively ample numerical divisor class.  By the BFJ identification
\(\operatorname{Div}(\pi)\simeq NS(Y/X)_{\mathbf R}\), this class has a
unique exceptional representative (A\), and (A\) is \(\pi\)-ample.  On
the vector space spanned by the (E_i\), define
\[
                         Q_A(D,C)=D\cdot C\cdot A^{n-2}.
\tag{12}
\]
We claim that (Q_A\) is negative definite.

To prove the claim, compactify \(\pi\) and resolve the boundary, without
changing a neighborhood of \(\pi^{-1}(0)\), obtaining a projective
birational morphism \(\bar\pi:\bar Y\to\bar X\) of smooth projective
(n\)-folds.  Choose a \(\bar\pi\)-ample Cartier divisor \(\bar A\), and
let (A\) be the exceptional representative of its relative numerical
class near the fibre over (0\), as above.  For a sufficiently ample divisor
(B\) on \(\bar X\),
\[
                         H=\bar A+\bar\pi^*B
\tag{13}
\]
is ample on \(\bar Y\).  On cycles supported in the exceptional fibre,
\(\bar A\) and (A\) have the same numerical class.  Every factor
\(\bar\pi^*B\) has zero intersection with such a cycle because the pulled
back line bundle restricts trivially to a fibre over a point.  Consequently,
for exceptional divisors (D,C\),
\[
 D\cdot C\cdot H^{n-2}=D\cdot C\cdot A^{n-2}=Q_A(D,C).
\tag{14}
\]

The generalized Hodge index theorem on the smooth projective variety
\(\bar Y\) says that
\[
 q_H(L_1,L_2)=L_1\cdot L_2\cdot H^{n-2}
\tag{15}
\]
has signature \((1,\rho(\bar Y)-1)\) on
\(N^1(\bar Y)_{\mathbf R}\).  The class \(\bar\pi^*B\) has positive square
for this form.  Each exceptional divisor (D\) is orthogonal to
\(\bar\pi^*B\), since the latter restricts trivially to (D\), which maps to
the point (0\).  Thus (15) is negative definite on the numerical span of
the exceptional divisors.  BFJ Section 1.2 identifies the coefficient space
of exceptional divisors with the relative numerical space, so no nonzero
exceptional divisor has zero numerical class.  It follows from (14) that
\(Q_A\) is negative definite, proving the claim.  As a check, for the blowup
of a smooth point, (A=-E\) and (E^2(-E)^{n-2}=-1\).

We finish the proof.  If (n=1\), the normalized valuation space consists
of the single valuation \(\operatorname{ord}_0\), and the generalized
Lelong number is the ordinary Lelong number of the tested germ; the result
is immediate.  Let (n\ge2\), fix any BFJ modification \(\pi:Y\to X\), and
write
\[
 D=Z(u)_\pi-Z(v)_\pi=\sum_{i=1}^r d_iE_i.
\tag{16}
\]
For every exceptional real divisor (C\) on (Y\), equation (11), applied
to the Cartier b-divisors determined by (C,A\), gives
\[
 \begin{aligned}
 0&=\Lambda(C,A,\ldots,A)\\
  &=\bigl(Z(u)_\pi-Z(v)_\pi\bigr)\cdot C\cdot A^{n-2}
   =Q_A(D,C).
 \end{aligned}
\tag{17}
\]
The nondegeneracy of (Q_A\) implies (D=0\).  Since \(\pi\) was arbitrary,
the traces of (Z(u)\) and (Z(v)\) agree on every modification; equivalently
their coefficients agree at every divisorial valuation.  Divisorial
valuations are dense in each quasimonomial simplex and formal psh functions
are continuous on each such simplex, so \(\widehat u\) and \(\widehat v\)
agree on the quasimonomial locus.  Finally BFJ extend a formal psh function
to every normalized valuation by
\[
                          g(\nu)=\inf_\pi g(r_\pi(\nu)).
\tag{18}
\]
The two extensions in (18) are equal.  Therefore
\(\widehat u=\widehat v\) pointwise on \(\mathcal V\).

The argument used only the analytic primary-ideal weights (3), proving the
last assertion as well.  Equality is pointwise, not merely almost everywhere
or quasi-everywhere.  The conclusion concerns valuative transforms; it does
not assert that (u-v\) is bounded, and bounded perturbations are invisible
on both sides from the outset.
